\question
В преддверии строительства инновационного города Токио-3, вы устраиваетесь на работу городского планировщика. Вам будет необходимо проектировать расположение домов и улиц будущего мегаполиса, то есть работать в первую очередь с графами. Чтобы оценить ваши архитектурные навыки, работодатель решил предложить вам тестовое задание: спроектировать и нарисовать следующие графы, а также обосновать корректность вашего ответа. Если же построить такой граф невозможно, необходимо строго обосновать, почему проект улиц и домов противоречит законам математики и не может быть реализован даже в условиях передовых технологий Токио-3.

\begin{enumerate}
\item  дерево с диаметром равным 6, двумя центрами  и  18-ю листьями;
\item  ориентированный граф с 5 компонентами слабой связности: в  каждой компоненте слабой связности не менее 5 компонент сильной связности;
\item  графы $K_6$ и $K_{6,6}$, $K_{3,4}$;
\item  простой неориентированный с 3 компонентами связности, который имеет 3 висячих ребра, 2 простых цикла длиной 7, 2 изолированные вершины, 3 моста и 8 точек сочленения.
\end{enumerate}

---------------

Автор -- Тимур Гонтарь
